\section{Introduction}\label{sec:introduction}

For an integer $n\geq2$, let $\Pplus(n)$ denote its largest prime divisor,
and put $\Pplus(1)=1$.  The Dickman--de Bruijn function $\rho$ is the
continuous function on $[0,\infty)$ determined by
\[
 \rho(u)=1\quad(0\leq u\leq1),\qquad
 u\rho'(u)=-\rho(u-1)\quad(u>1).
\]
The classical theory of smooth
numbers, beginning with Dickman and developed by Ramaswami and
de Bruijn~\cite{Dickman1930,Ramaswami1949,deBruijn1951}, gives
\[
 \frac1X\#\{n\leq X:\Pplus(n)\leq X^a\}
 \longrightarrow \rho(1/a)\qquad(0<a\leq1).
\]
Our result identifies the joint law at two consecutive integers.

\begin{theorem}[Joint Dickman law]\label{thm:main}
For every fixed $a,b\in(0,1)$,
\[
 \lim_{X\to\infty}\frac1X
 \#\{2\leq n\leq X:\Pplus(n)\leq n^a,
                         \ \Pplus(n+1)\leq n^b\}
 =\rho(1/a)\rho(1/b).
\]
Here the limit is taken over all real $X$, with ordinary, unweighted
counting.
\end{theorem}

Thus $\log\Pplus(n)/\log n$ and $\log\Pplus(n+1)/\log n$ have, in natural
density, independent limiting distributions with distribution function
$a\mapsto\rho(1/a)$, continuously extended to $[0,1]$.
Section~\ref{sec:distribution} proves the corresponding law with fixed
thresholds $X^a,X^b$, derives the displayed moving-threshold statement,
and identifies this continuous limiting law.  Inclusion--exclusion using
the fixed-threshold law and its marginals gives the upper-tail independence conjecture of
Erd\H{o}s and Pomerance~\cite[p.~311]{ErdosPomerance1978}.

The comparison conjecture usually attributed to Erd\H{o}s and Tur\'an is
a consequence, rather than a separate main theorem.

\begin{corollary}\label{cor:comparison}
One has
\[
 \lim_{X\to\infty}\frac1X
 \#\{2\leq n\leq X:\Pplus(n)<\Pplus(n+1)\}=\frac12.
\]
The reverse ordering has the same natural density.
\end{corollary}

The limiting marginal is continuous, so its product measure gives zero
mass to the diagonal; symmetry then gives the corollary.  In particular,
no quantitative separation estimate is needed for this deduction.

\subsection{Previous work}

Erd\H{o}s and Pomerance~\cite{ErdosPomerance1978} formulated the joint
independence problem and proved that each ordering has positive lower
natural density.  Their explicit lower bound was $0.0099$.  They also
proved that for every $\varepsilon>0$, some $\delta>0$ makes the number of
integers $n<X$ satisfying
$X^{-\delta}<\Pplus(n)/\Pplus(n+1)<X^\delta$ less than $\varepsilon X$
for all sufficiently large $X$~\cite[Theorem~1]{ErdosPomerance1978}.
The lower-density bound was increased to $0.05544$ by de la Bret\`eche,
Pomerance and Tenenbaum~\cite[Section~3]{BretechePomeranceTenenbaum2005},
and to $0.05866$ by an observation of Fouvry recorded in the same section.
Wang~\cite{Wang2017,Wang2018} obtained $0.1063$ and $0.1356$, and
L\"u and Wang~\cite{LuWang2025} obtained $0.2017$.
Yang's recent preprint~\cite[Theorem~1.4]{Yang2026} gives $0.280$ for both
orderings.  These are bounds on lower natural densities; they do not
assert that a natural density exists.

For the joint law itself,
Ter\"av\"ainen~\cite[Theorem~1.14]{Teravainen2018} proved the predicted
product in logarithmic density.  That is, for fixed $0<a,b<1$, the
indicator in Theorem~\ref{thm:main}, averaged with weight $1/n$ and
normalization $1/\log X$, tends to $\rho(1/a)\rho(1/b)$.
His Theorem~1.16 gives logarithmic density $1/2$ for the ordering, while
Theorem~1.19 proves positive lower natural density for every nondegenerate
rectangle of normalized largest-prime-factor values inside $(0,1)^2$.
Tao and Ter\"av\"ainen~\cite[Remark~3.3, equation~(50)]{TaoTeravainen2019}
then obtained the joint product law for ordinary averages outside an
exceptional set of scales of logarithmic density zero; their
Corollary~1.16 gives the corresponding ordering result.
Wang~\cite{Wang2021Conditional} proved the ordinary joint law under the
Elliott--Halberstam conjecture for friable integers.  Jiang, L\"u and
Wang~\cite{JiangLuWang2022} established averaged-over-shift forms of the
conjectures, a different conclusion from a result at the fixed shift one.

The more recent work of Tao and
Ter\"av\"ainen~\cite[Theorem~1.8]{TaoTeravainen2026Quantitative} gives
a quantitative joint law outside an exceptional set of scales, with
uniformity in the smoothness parameters and a power saving in $\log X$.
Theorem~\ref{thm:main} concerns fixed parameters and has no exceptional
scales; it does not assert their quantitative error term or their
uniformity for growing parameters.  Relative to the results just
described, the distinction is therefore the unconditional ordinary limit
at every sufficiently large scale, not a new marginal distribution or
a logarithmically averaged independence statement.

\subsection{Method and organization}

We first replace the largest prime factor by finitely many prime counts.
For a fixed integer $J\ge2$, divide the primes in $(x^{1/J},x]$ into
bins $\mathcal B_{k,x}=(x^{k/J},x^{(k+1)/J}]$, $1\le k<J$.  If
$\Omega_E(n)$ counts prime factors in $E$ with multiplicity, a choice of
unit complex numbers $\zeta_k$ defines the bin character
\[
 f_x(n)=\prod_{k=1}^{J-1}\zeta_k^{\Omega_{\mathcal B_{k,x}}(n)}.
\]
Let $g_x$ be defined by a second arbitrary phase vector, and
let $F_x=f_x-\mu$, where $\mu$ is the limiting mean of $f_x$.  Finite
Fourier inversion reduces the joint law to the mixed decorrelation
\[
 \frac1x\sum_{n<x}\overline{g_x(n)}F_x(n+1)\longrightarrow0.
\]
Here $x$ runs through integer scales.
The two phase vectors may differ.  This freedom gives independence of
the two count vectors.

Two properties of these functions connect the beginning and end of the
proof.  First, every fixed positive integer multiplier has all its prime factors
below the bins once $x$ is large.  Thus $F_x(un)=F_x(n)$ and
$g_x(un)=g_x(n)$ for each fixed $u$.  Second, $F_x$ has cancellation in
averages over growing intervals that remain short relative to the original
scale.  Let $B$ be the auxiliary parameter and let $T=T(B)\to\infty$ be
the growing shift scale.  If $P_B$ is a finite set of primes with
$\min P_B\to\infty$, define
\[
 G_{B,v}(n)=\prod_{\substack{p\in P_B\\p\mid n}}
                 \frac{1+p^{-v/B}}2,\qquad v\ge0.
\]
For any interval $I_B$ of $L_B\to\infty$ consecutive integer shifts,
and fixed $0<s_0<s_1$, positive integer $D$, residue $r\pmod D$, and $v\ge0$,
Section~\ref{sec:labels} proves
\[
 \lim_{B\to\infty}\limsup_{x\to\infty}
 \frac1{Tx}\sum_{s_0Tx\le n\le s_1Tx}
 \left|\frac1{L_B}
   \sum_{\substack{i\in I_B\\i\equiv r\pmod D}}
     F_x(n+i)G_{B,v}(n+i)\right|^2=0.
\]
At fixed $B$ the shift interval is short relative to $x$, while the base
points have size $Tx$.  This is the scale produced by the divisor
substitution below.

The centered function $F_x$ is not multiplicative.  On the relevant ranges
$n=O_B(x)$, the bin counts take values in a fixed finite set.  This permits
interpolation by a fixed finite family of real,
nonnegative multiplicative functions.  After resolving the residue
condition by Dirichlet characters, the real short-interval theorem of
Matom\"aki and Radziwi\l\l~\cite[Theorem~1]{MatomakiRadziwill2016} compares
the principal components with long means whose centered combination
vanishes.  The complex theorem of Matom\"aki, Radziwi\l\l\ and
Tao~\cite[Theorem~A.1, corrected version]{MatomakiRadziwillTao2015}
controls the nonprincipal twists.  The rest of the proof must reduce
the mixed correlation to this particular cancellation statement.

Suppose, then, that the mixed correlation stays nonzero along a sequence
of scales.  The profinite integers $\widehat{\mathbb Z}$ encode compatible
residue classes modulo all positive integers and carry Haar probability
measure.  Passing to a subsequence gives a bounded profile $W(t,w)$ on
$(0,\infty)\times\widehat{\mathbb Z}$.  Its integral against each fixed
compactly supported continuous test $\Phi(t,w)$ is the subsequential limit
of the mixed averages weighted by $\Phi(n/x,n)$, with $n$ in the second
argument viewed through its residues.  A fixed nonnegative compactly
supported smooth cutoff $\phi$ can be chosen so that
$\int\phi(t)W(t,w)\,dt\,dw\ne0$, where $dt$ is Lebesgue measure and $dw$
is Haar measure.

For an integer $n$, the amplifier defines $D_B(n)$ as a nonnegative
weighted count of factorizations $n=am$ and $n+1=cl_a$, restricted by
$e^B<c<e^{2B}$ and $Tc<a<2Tc$.  At fixed $B$ its coefficient list is
finite, and the same formula defines a function on $\widehat{\mathbb Z}$
depending only on residues at primes above a cutoff tending to infinity.
The construction gives $D_B$ Haar mean at least $d_0>0$ for large $B$ and
bounded $L^2$ norm.  To prove these bounds, the argument passes to a
comparison model in which each site prime is assigned by an independent
fair coin to the coefficient divisor or the remaining factor.  The first
and second moments are analyzed through one split and two conditionally
independent splits of the same site prime sets.  For each fixed $B$,
$\Phi(t,w)=\phi(t)D_B(w)$ is an allowed profile test, so along the selected
subsequence
\[
 \lim_{x\to\infty}\frac1x\sum_{n\ge1}
   \phi(n/x)\overline{g_x(n)}F_x(n+1)D_B(n)
 =\int\phi(t)W(t,w)D_B(w)\,dt\,dw.
\]
Independence from every fixed residue coordinate and $L^2$ approximation
of $\phi W$ show that the right-hand side has absolute value bounded away
from zero as $B\to\infty$.

For fixed $B$ and large $x$, writing $n=am$ in this weighted average
replaces $g_x(n)$ by $g_x(m)$.
For each fixed pair $(c,m)$, the remaining terms form a sum over
$a$ satisfying $c\mid am+1$, with nonnegative divisor and smooth weights
multiplying $F_x(am+1)$.  Cauchy--Schwarz in $(c,m)$ removes the
unit-modulus factor $g_x(m)$ and squares this inner sum.  The resulting
normalized quadratic energy has a positive lower bound, while its
equal-coefficient diagonal tends to zero in the same iterated limit.  A positive contribution
therefore remains from distinct indices $a,b$.
Section~\ref{sec:amplification} proves these moment and energy claims.

For such a pair, $c$ divides both $am+1$ and $bm+1$.  Thus $m$ is a unit
modulo $c$ and $a-b=jc$, with $0<|j|<T$ by the coefficient windows.  Writing
$l_a=(am+1)/c$ and $l_b=(bm+1)/c$, the exact substitution
\[
 n'=b l_a,\qquad n'+j=a l_b
\]
turns the label product into $\overline{F_x(n')}F_x(n'+j)$ by invariance
under the fixed multipliers $a,b,c$.  The new endpoints have size $Tx$.
The off-diagonal energy is therefore a weighted graph of ordinary additive
shifts.  Its edge weights depend on endpoint prime sets and on additional
residue data from each divisor representation.

Group the new endpoints into blocks $N+i$, $1\le i\le M$, with
$M=\lceil C_6T\rceil$ for a fixed large $C_6$, and put $s=N/(Tx)$.
Section~\ref{sec:graph} couples the actual auxiliary-prime divisibility
sets at these positions to independent model sets $S_i$, each
formed by including every auxiliary prime $p$ independently with probability $1/p$, and averages
the extra residue weight of each representation.  This compares the
arithmetic graph with an endpoint kernel
$\mathcal L_{ik}(S_i,S_k;s)$ in expected cut norm.  Here cut norm is the
largest absolute normalized pairing with two real vectors bounded by one,
chosen after the matrix is known.  Section~\ref{sec:moments} controls
representation multiplicity and provides the averaged second-moment
bounds needed for sampling.

The aim is now to replace this endpoint kernel by products of functions of
its two endpoints.  For an endpoint type $S$, let $b$ be the product
obtained by retaining each prime independently
with probability $1/2$.  The features used in the comparison are
\[
 V_l(S)=\frac1{|I_l|}
   \Prob_{\rm split}\!\left(\frac{\log b}{B}\in I_l\,\middle|\,S\right),
\]
where the $I_l$ form a fixed coarse partition of a compact logarithmic
interval, chosen for a desired accuracy.  The channel estimates in Section~\ref{sec:channels} show that
fair splitting suppresses residue dependence and makes logarithmic
outputs uniformly approximable on this partition.  Fourier detection of
$a-b=jc$ then lets Section~\ref{sec:kernel} apply arithmetic estimates to
the coefficient $c$ and these channel estimates to the endpoints.  The
resulting comparison has the form
\[
 \sum_l c_{l,B}(j,s)V_l(S_i)V_l(S_k),\qquad j=k-i,
\]
where $i,k$ are the endpoint positions.  The scalar coefficients separate
an arithmetic factor in $|j|$ from a smooth factor in the normalized lag
$|j|/T$ and the common block origin $s$.  The number of features is fixed,
while their values and the coefficients may depend on $B$.

The feature comparison first holds after integration against separate
functions of the two independent endpoint sets.  The actual labels need
not have this single-site form, and an optimizing test can depend on the
whole configuration.  Section~\ref{sec:sampling} upgrades the comparison
to expected cut norm by approximating an optimizing cut with a small sample
of columns.  Finally,
polynomial approximation of each log-window indicator expresses $V_l$,
up to a small mean-square error, as a fixed finite combination of
\[
 \E_{\rm split}\!\left[e^{-v\log b/B}\mid S\right]
       =\prod_{p\in S}\frac{1+p^{-v/B}}2.
\]
At an integer endpoint these are exactly $G_{B,v}$.  The arithmetic lag
factor is approximated in mean over lags by a fixed periodic function.
For a fixed small $\delta>0$ chosen for the desired accuracy, subdivide each
position block into a fixed number of intervals whose lengths grow with
$T$ and are at most $\delta T$.  Within each interval pair, the normalized
lag varies by at most $2\delta$.  Approximating the smooth lag factor there
and resolving the periodic factor into fixed residue classes turns the
feature energy, up to a controlled error, into products of the weighted
short averages already controlled in Section~\ref{sec:labels}.  Their vanishing
contradicts the positive energy.  Sections~\ref{sec:conclusion}
and~\ref{sec:distribution} complete this argument and the passage from bin
events to the joint law.

The divisor amplification and graph comparison have close antecedents in
Tao's logarithmic correlation argument~\cite{Tao2016LogChowla}, the
prime-divisibility graphs of Helfgott and
Radziwi\l\l~\cite{HelfgottRadziwill2021}, Pilatte's product-of-primes
amplification~\cite{Pilatte2023}, and the mixed-function decoupling of
Tao and Ter\"av\"ainen~\cite[Section~3.1]{TaoTeravainen2026Quantitative}.
The present rough-divisor graph and its endpoint comparison are proved
locally.  The sampling argument likewise builds on the cut-based methods
of Frieze and Kannan~\cite{FriezeKannan1999}, Alon, Fernandez de la Vega,
Kannan and Karpinski~\cite[Lemma~3]{AlonEtAl2003Sampling}, and Borgs,
Chayes, Lov\'asz, S\'os and
Vesztergombi~\cite[Theorem~4.6]{BorgsEtAl2008}.  Because the kernels here
vary with $B$ and can be large, the proof establishes the moment bounds
needed before applying bounded-differences concentration~\cite{McDiarmid1989}.

The arithmetic estimates are developed in Section~\ref{sec:arithmetic}
using Selberg--Delange theory, character estimates and upper-bound sieves
in the forms recorded by Granville and
Koukoulopoulos~\cite{GranvilleKoukoulopoulos2019},
Koukoulopoulos~\cite{Koukoulopoulos2019}, and
Ford~\cite{FordSieve2023,FordZeta2022}.  The minor-arc input in
Section~\ref{sec:kernel} is the exponential-sum estimate of Montgomery and
Vaughan~\cite{MontgomeryVaughan1977}.  Throughout the proof, the original
scale $x$ tends to infinity before the auxiliary scale $B$.  Bin data,
feature degrees, residue moduli and divisor multipliers are fixed in that
inner limit.  The contradiction applies to a subsequence of every possible
bad sequence of original scales, and hence yields the full ordinary limit.
